\section{Assumptions and Dense-Span Property}
\label{app:assumptions}

This section records the assumptions used by the convergence theorem and fixes the notation for the mean-field limit.

We collect the assumptions used in Theorem~\ref{thm:convergence}. We reuse the multilayer mean-field framework and notation of Nguyen and Pham~\cite{nguyen2023}: the mean-field width limit sends $n_1,\ldots,n_{L-1}\to\infty$ at fixed depth $L$, and discrete neuron indices are replaced by labels $C_\ell\in\Omega_\ell$ with laws $\rho_\ell$. Our only architectural change is that dense contractions are replaced by frozen random features and bounded low-rank channel mixers. The notation $\partial_2\calL(y,u)$ denotes the derivative of the loss with respect to its prediction argument.

\begin{assumption}[Bounded activations and low-rank mixing]
\label{assump:regularity}
There exists $K\ge1$ such that the activations $\varphi_\ell$ are $K$-Lipschitz, the activation derivatives used in backpropagation are bounded on the active region, and the low-rank mixers satisfy
\begin{equation}
\sup_{c_\ell}\sum_{k=1}^{r_\ell}|L^{(\ell)}_{c_\ell,k}|\le r_\ell K .
\end{equation}
The mixers are full-column-rank tall-skinny factors in finite width. Orthogonality, $L^{(\ell)\top}L^{(\ell)}=I_{r_\ell}$, is an optional Stiefel normalization rather than a theorem assumption. The preactivations remain in the a priori bounded regime on every finite interval, and the learning-rate schedules are non-negative, bounded, and locally integrable.
\end{assumption}

\begin{assumption}[Sub-Gaussian initialization]
\label{assump:init}
The trainable initial weights have sub-Gaussian tails, uniformly over channels. In the two-hidden-layer notation,
\begin{equation}
\sup_{m\ge1}\frac{1}{\sqrt m}
\max_{1\le k\le r}
\E_{C_1}\!\left[|w_1^0(C_1,k)|^m\right]^{1/m}\le K,
\qquad
\sup_{m\ge1}\frac{1}{\sqrt m}
\E_{C_2}\!\left[|w_2^0(C_2)|^m\right]^{1/m}\le K .
\end{equation}
The $L$-layer version imposes the same moment growth bound on every trainable channel.
This assumption includes the standard initializations used in practice. In particular, Xavier or Glorot uniform initialization~\cite{glorot2010}, with entries sampled independently from
\[
\operatorname{Unif}\!\left[-\sqrt{\frac{6}{d_{\rm in}+d_{\rm out}}},
\sqrt{\frac{6}{d_{\rm in}+d_{\rm out}}}\right],
\]
is bounded and therefore sub-Gaussian after the usual fan-in/fan-out scaling. Xavier normal initialization, and more generally any independent centered sub-Gaussian initialization with the same variance scale, also satisfies the displayed moment bound. Thus the phrase ``standard independent initialization'' in Theorem~\ref{thm:convergence} covers Glorot/Xavier uniform, Glorot/Xavier normal, and the usual sub-Gaussian or subnormal variants used for trainable weights.
\end{assumption}

\begin{assumption}[Data distribution and loss]
\label{assump:data}
\label{assump:loss}
Inputs are bounded, $|X|\le K$ almost surely, and the frozen first-layer features satisfy $\|L^0(c_1)\|\le K$. The loss is non-negative, and $\partial_2\calL(y,\cdot)$ is bounded and Lipschitz on the relevant prediction range. Moreover,
\begin{equation}
\E[\partial_2\calL(Y,u)\mid X=x]=0
\quad\Longrightarrow\quad
\E[\calL(Y,u)\mid X=x]=0
\label{eq:loss-identifiability}
\end{equation}
for $\Pp_X$-almost every $x$.
This condition is the usual realizable or excess-risk identifiability condition. For squared loss, it holds in the noiseless realizable case, where $Y=f^\star(X)$ and $\partial_2 (Y-u)^2=0$ implies $u=Y$ and therefore zero conditional loss; equivalently, in the noisy case it holds after replacing the raw squared loss by its excess risk above the Bayes regressor. For cross-entropy, the same statement holds when the loss is written as the excess cross-entropy, or conditional KL divergence, relative to the Bayes conditional label distribution; in the deterministic-label realizable case this again reduces to zero conditional loss at the correct classifier.
\end{assumption}

\begin{assumption}[Diversity of frozen random features]
\label{assump:diversity}
The support of the law of $L^0(C_1)$ is dense in $\R^d$, and $\varphi_1$ is non-polynomial. Equivalently,
\begin{equation}
\left\{
\varphi_1\!\left(\langle L^0(c_1),\cdot\rangle\right):c_1\in\Omega_1
\right\}
\end{equation}
has dense span in $L^2(\Pp_X)$.
\end{assumption}

\begin{assumption}[Non-degenerate limit]
\label{assump:lr}
The limiting representation is not a dead network:
\begin{equation}
\max_{1\le k\le r}\Pp(\bar w_1(C_1,k)\ne0)>0,
\qquad
\Pp(\bar w_\ell(C_\ell)\ne0)>0,\quad \ell=2,\ldots,L-1 .
\end{equation}
A sufficient condition is that the initial loss is strictly better than the trivial predictor.
\end{assumption}

\begin{assumption}[Convergence in a modified $\mathcal W_4$ coupling topology]
\label{assump:convergence}
The mean-field trajectory has a limit $\bar W$ in a modified $\mathcal W_4$ topology. This is a Wasserstein-4-type coupling topology, with the fourth-moment control replaced by the weighted channel quantities that appear in the low-rank ODE and in the output-stability estimates. In the two-hidden-layer notation, there exist couplings $\pi_t$ such that
\begin{align}
\int &(1+|\bar w_2(c_2)|)|\bar w_2(c_2)|
\max_{1\le k\le r}|\bar w_1(c_1,k)|
|w_1(t,c_1',k)-\bar w_1(c_1,k)|\,d\pi_t
\to0,
\label{eq:conv-w1-neurips}\\
\int &(1+|\bar w_2(c_2)|)|\bar w_2(c_2)|
|w_2(t,c_2')-\bar w_2(c_2)|\,d\pi_t
\to0 .
\label{eq:conv-w2-neurips}
\end{align}
For depth $L>3$, the assumption includes the analogous weighted coupling gaps for all trainable layers.
\end{assumption}

\begin{theorem}[Universal approximation automatically maintained]
\label{thm:uap-automatic}
Under Assumption~\ref{assump:diversity}, the frozen first-layer class has dense span in $L^2(\Pp_X)$ throughout training.
\end{theorem}

\begin{proof}
This is the classical non-polynomial random-feature density theorem. Since the first feature map is frozen, its support cannot collapse during training.
\end{proof}

\section{Supplementary Proof of Well-Posedness}
\label{app:wellposed-proof}

This section gives the full Picard fixed-point proof that the low-rank mean-field ODE is well defined on every finite time interval.

We give the complete argument for the existence and uniqueness part of Theorem~\ref{thm:convergence}. The proof is deliberately close to the Picard proof for multilayer mean-field networks in Nguyen and Pham~\cite{nguyen2023}. The low-rank architecture does not introduce a new analytic difficulty: every place where the full-rank proof uses a dense-layer operator norm, we use the bounded row-sum constant
\begin{equation}
\Lambda_\ell
:=
\sup_{c_\ell}\sum_{k=1}^{r_\ell}|L^{(\ell)}_{c_\ell,k}|
\le r_\ell K .
\label{eq:lr-mixing-constant}
\end{equation}
Thus the proof is a natural specialization of the existing mean-field Picard argument to bounded low-rank mixing.

For clarity we write the proof in the two-hidden-layer notation. The extension to arbitrary fixed depth is obtained by repeating the same estimates layer by layer and replacing $\Lambda_2$ by $\max_\ell \Lambda_\ell$ in the constants. Recall
\begin{equation}
f_k(t,x)=
\E_{C_1}\!\left[
w_1(t,C_1,k)\,
\varphi_1(\langle L^0(C_1),x\rangle)
\right],
\qquad
H_2(t,c_2;x)=\sum_{k=1}^r L_{c_2,k}f_k(t,x),
\end{equation}
and
\begin{equation}
\hat y(x;W(t))
=
\E_{C_2}\!\left[
w_2(t,C_2)\varphi_2(H_2(t,C_2;x))
\right].
\end{equation}
Let $G(W)$ denote the right-hand side of the mean-field ODE. In integral form, a solution is a fixed point of
\begin{equation}
F(W)(t)=W(0)+\int_0^t G(W(s))\,ds .
\label{eq:picard-map}
\end{equation}

\paragraph{Step 1: low-rank forward Lipschitz estimates.}
Let $W'$ and $W''$ be two trajectories on $[0,T]$. Since $\varphi_1$ is bounded on the relevant input range and $|X|,\|L^0(C_1)\|\le K$,
\begin{equation}
|f'_k(t,x)-f''_k(t,x)|
\le
K\,\E_{C_1}|w'_1(t,C_1,k)-w''_1(t,C_1,k)|.
\end{equation}
Consequently the low-rank preactivation satisfies
\begin{align}
|H'_2(t,c_2;x)-H''_2(t,c_2;x)|
&\le
\sum_{k=1}^r |L_{c_2,k}|\,
|f'_k(t,x)-f''_k(t,x)|
\nonumber\\
&\le
K\Lambda_2
\max_{1\le k\le r}
\E_{C_1}|w'_1(t,C_1,k)-w''_1(t,C_1,k)|.
\label{eq:H-lr-lipschitz}
\end{align}
This is the only low-rank modification of the dense proof. The dense operator norm is simply replaced by $\Lambda_2$.

\paragraph{Step 2: drift Lipschitz estimates on bounded sets.}
Fix $T<\infty$ and consider the class $\mathcal W_T^0$ of trajectories with the same initialization, bounded fourth moments, and sub-Gaussian tails on $[0,T]$. More explicitly, there is a finite constant $K_0(T)$ such that, for every $W\in\mathcal W_T^0$,
\begin{equation}
\sup_{t\le T}
\left(
\max_k \E_{C_1}|w_1(t,C_1,k)|^4
+
\E_{C_2}|w_2(t,C_2)|^4
\right)^{1/4}
\le K_0(T),
\end{equation}
with the corresponding uniform sub-Gaussian tail bound inherited from the initialization and the bounded drift. On this class, the loss derivative, activations, activation derivatives, and low-rank mixers are bounded or Lipschitz by Assumptions~\ref{assump:regularity}--\ref{assump:data}. Combining these bounds with \eqref{eq:H-lr-lipschitz} gives, on the good event where the channel weights are bounded by $K_0(T)B$,
\begin{equation}
\|G(W')-G(W'')\|_t
\le
C_T(1+B)\|W'-W''\|_t .
\label{eq:drift-good}
\end{equation}
Here $C_T$ depends on $K$, $T$, and the low-rank constants $\Lambda_\ell$, but not on width. The complement of the good event is controlled by the sub-Gaussian tail bound, hence contributes at most $C_T e^{-cB^2}$. Therefore
\begin{equation}
\|F(W')-F(W'')\|_t
\le
C_T(1+B)\int_0^t \|W'-W''\|_s\,ds
+
C_T e^{-cB^2}.
\label{eq:picard-difference}
\end{equation}
This is the same estimate used in Nguyen and Pham~\cite{nguyen2023}; only the constant $C_T$ changes through \eqref{eq:lr-mixing-constant}.

\paragraph{Step 3: invariance of the Picard map.}
The map $F$ sends $\mathcal W_T^0$ into itself. Indeed, the integral formula \eqref{eq:picard-map}, the boundedness of $\partial_2\calL$, the boundedness of the activations on the a priori regime, and the low-rank estimate
\begin{equation}
|H_2(t,c_2;x)|
\le
\Lambda_2 \max_{1\le k\le r}|f_k(t,x)|
\le
K\Lambda_2
\max_{1\le k\le r}\E_{C_1}|w_1(t,C_1,k)|
\end{equation}
imply a Gronwall bound for the fourth moments on $[0,T]$. The same argument applied to exponential moments gives the sub-Gaussian tail bound. Thus, after increasing $K_0(T)$ if necessary, $F(W)\in\mathcal W_T^0$ whenever $W\in\mathcal W_T^0$.

\paragraph{Step 4: Picard convergence.}
Iterating \eqref{eq:picard-difference} yields, for all $m\ge1$,
\begin{align}
\|F^m(W')-F^m(W'')\|_T
&\le
\frac{(C_TT(1+B))^m}{m!}\|W'-W''\|_T
+
C_T\exp(C_TT(1+B)-cB^2).
\end{align}
Choose $B=\sqrt m$. The first term tends to zero by Stirling's formula and the second tends to zero exponentially. Taking $W''=F(W')$, the series
\begin{equation}
\sum_{m\ge1}\|F^{m+1}(W')-F^m(W')\|_T
\end{equation}
is finite, so the Picard iterates converge in $\|\cdot\|_T$ to a limit $W$. Since $F$ is continuous under the same estimate, $F(W)=W$, hence $W$ solves the mean-field ODE on $[0,T]$.

\paragraph{Step 5: uniqueness and continuation.}
If $W'$ and $W''$ are two fixed points of $F$ in $\mathcal W_T^0$, then applying the iterated estimate to the fixed points gives
\begin{equation}
\|W'-W''\|_T
=
\|F^m(W')-F^m(W'')\|_T
\to0
\qquad (m\to\infty).
\end{equation}
Thus the solution is unique on $[0,T]$. Since $T<\infty$ was arbitrary and the a priori bounds are finite on every finite interval, the solution extends uniquely to $[0,\infty)$. This proves the well-posedness part of Theorem~\ref{thm:convergence}.

\section{Proof Details for Global Convergence}
\label{app:global-proof}

This section expands the stationarity, dense-span, non-degeneracy, and coupling arguments behind the global convergence theorem. We write the proof in the two-hidden-layer notation to keep the formulas readable. The same argument extends to any fixed depth because the decisive step is the frozen first layer: the class $\{\varphi_1(\langle L^0(c_1),\cdot\rangle):c_1\in\Omega_1\}$ keeps dense span throughout training, so the stationarity identity can always be pushed back to this unchanged feature family. Additional hidden layers only add backpropagated scalar factors and low-rank mixing constants; they do not change the dense-span-to-optimality mechanism.

\begin{proof}[Proof of Theorem~\ref{thm:convergence}]
Let $Z=(X,Y)$ and write
\begin{equation}
d_L(Z;W)=\partial_2\calL(Y,\hat y(X;W)).
\end{equation}
In the two-hidden-layer notation,
\begin{equation}
f_k(t,x)=
\E_{C_1}\!\left[
w_1(t,C_1,k)\,
\varphi_1(\langle L^0(C_1),x\rangle)
\right],
\qquad
H_2(t,c_2;x)=\sum_{k=1}^r L_{c_2,k}f_k(t,x),
\end{equation}
and
\begin{equation}
\hat y(x;W(t))=
\E_{C_2}\!\left[
w_2(t,C_2)\,
\varphi_2(H_2(t,C_2;x))
\right].
\end{equation}

The well-posedness part is proved in Appendix~\ref{app:wellposed-proof}. It is exactly the natural Picard argument of Nguyen and Pham~\cite{nguyen2023}, with dense-layer operator norms replaced by the bounded low-rank mixing constants. The dense-span property is preserved because $L^0(C_1)$ is never trained.

Let $W(t)\to\bar W$ in the modified $\mathcal W_4$ coupling topology of Assumption~\ref{assump:convergence}. At the limit, stationarity of the first trainable low-rank channel gives, for every $c_1$ and $k$,
\begin{equation}
\E_Z\!\left[
d_L(Z;\bar W)\,
\varphi_1(\langle L^0(c_1),X\rangle)\,
B_k^{(2)}(X;\bar W)
\right]=0,
\label{eq:app-first-stationary}
\end{equation}
where
\begin{equation}
B_k^{(2)}(x;\bar W)
:=
\E_{C_2}\!\left[
L_{C_2,k}\,
\varphi_2'(H_2(C_2;x,\bar W))\,
\bar w_2(C_2)
\right].
\label{eq:app-Bk}
\end{equation}
For deeper networks $B_k^{(\ell)}$ is the analogous backpropagated scalar. Conditioning on $X$ and using the dense-span property implies
\begin{equation}
\E\!\left[
d_L(Z;\bar W)\,
B_k^{(2)}(X;\bar W)
\mid X=x
\right]=0
\quad\text{for }\Pp_X\text{-a.e. }x,\quad k=1,\ldots,r .
\label{eq:app-dense-product}
\end{equation}
The conditional identity is kept in the form
\begin{equation}
\E\!\left[
d_L(Z;\bar W)\,
B_k^{(2)}(X;\bar W)
\mid X=x
\right]=0
\quad\text{for }\Pp_X\text{-a.e. }x,\quad k=1,\ldots,r .
\label{eq:app-dense-product-restated}
\end{equation}
Since $B_k^{(2)}(X;\bar W)$ is $X$-measurable and the non-degeneracy assumption gives at least one nonzero backpropagated factor almost everywhere, \eqref{eq:app-dense-product-restated} implies
\begin{equation}
\E\!\left[
\partial_2\calL(Y,\hat y(X;\bar W))
\mid X=x
\right]=0
\quad\text{for }\Pp_X\text{-a.e. }x .
\label{eq:app-conditional-zero}
\end{equation}
Assumption~\ref{assump:loss}, through the exact loss-identifiability implication \eqref{eq:loss-identifiability}, then gives $\E[\calL(Y,\hat y(X;\bar W))\mid X=x]=0$, so $\scrL(\bar W)=0$. Since $\calL\ge0$, $\bar W$ is a global minimizer.

It remains to connect the trajectory to the limit. A representative coupling gap is
\begin{equation}
\E_{\pi_t}\!\left[
(1+|\bar w_2|)\,|\bar w_2|\,
\sum_{k=1}^{r}|\bar w_{1,k}|\,|w_{1,k}(t)-\bar w_{1,k}|
\right]\longrightarrow0.
\label{eq:channel-coupling-gap}
\end{equation}
The low-rank form gives
\begin{equation}
H_2(c_2;x,W(t))-H_2(c_2;x,\bar W)
=
\sum_{k=1}^r L_{c_2,k}
\left(f_k(t,x)-\bar f_k(x)\right).
\end{equation}
Using Lipschitz activations and iterating through layers,
\begin{equation}
\E_Z\!\left[
|\hat y(X;W(t))-\hat y(X;\bar W)|
\right]\le K\,\Gamma_t\longrightarrow0,
\label{eq:output-coupling-gap}
\end{equation}
where $\Gamma_t$ is the sum of the weighted coupling gaps. Lipschitzness of the loss yields $\scrL(W(t))\to\scrL(\bar W)$.
\end{proof}

\section{Detailed Proof of the CPWL Fourier Shell Law}
\label{app:cpwl-proof}

This section proves the Fourier shell law by integrating by parts over the face lattice of a windowed CPWL function.

We prove Proposition~\ref{prop:shell-law}. Write the CPWL network on its polyhedral complex as
\begin{equation}
f(x)=\sum_{\varepsilon}
(a_\varepsilon^\top x+b_\varepsilon)\mathbf 1_{P_\varepsilon}(x),
\qquad
g(x)=\chi(x)f(x).
\end{equation}
The cutoff makes $g$ compactly supported, so
\begin{equation}
\wh g(\rho\theta)
=
\sum_\varepsilon
\int_{P_\varepsilon}
e^{-i\rho\langle \theta,x\rangle}
\chi(x)(a_\varepsilon^\top x+b_\varepsilon)\,dx
\end{equation}
is an ordinary oscillatory integral. Set $D_\theta=\theta\cdot\nabla$. Since
$D_\theta e^{-i\rho\langle\theta,x\rangle}=-i\rho e^{-i\rho\langle\theta,x\rangle}$,
the divergence theorem on each cell gives
\begin{align}
\int_{P_\varepsilon}e^{-i\rho\langle\theta,x\rangle}u_\varepsilon(x)\,dx
&=
\frac{i}{\rho}
\int_{\partial P_\varepsilon}
e^{-i\rho\langle\theta,x\rangle}
u_\varepsilon(x)\langle\theta,n_\varepsilon(x)\rangle\,d\sigma(x)
\nonumber\\
&\quad
-
\frac{i}{\rho}
\int_{P_\varepsilon}
e^{-i\rho\langle\theta,x\rangle}
D_\theta u_\varepsilon(x)\,dx,
\label{eq:first-ibp-cell-app}
\end{align}
where $u_\varepsilon(x)=\chi(x)(a_\varepsilon^\top x+b_\varepsilon)$ and $n_\varepsilon$ is the outward normal. When the boundary terms are summed over all cells, every interior facet $F=P_\varepsilon\cap P_{\varepsilon'}$ is counted twice with opposite normals. The order-$\rho^{-1}$ trace contribution is proportional to
\begin{equation}
\left(u_\varepsilon|_F-u_{\varepsilon'}|_F\right)
\langle \theta,n_F\rangle .
\end{equation}
Because $f$ is continuous and the same smooth window $\chi$ multiplies both traces, this term vanishes on internal facets. Exterior window terms are smooth cutoff terms and are absorbed in the final remainder.

The leading non-smooth contribution comes from applying the same identity to the integral involving $D_\theta u_\varepsilon$. On a facet $F$ shared by two cells, its jump is
\begin{equation}
\big[D_\theta u\big]_F
=
\chi\,\theta^\top(a_\varepsilon-a_{\varepsilon'})
\end{equation}
because the derivatives of $\chi$ multiply the continuous trace of $f$ and cancel across $F$. Thus the first non-canceling boundary term is a jump of directional derivatives and carries the factor $\rho^{-2}$.

Higher-codimension terms follow by induction on the face lattice. Assume that after $m$ reductions the surviving terms are sums over codimension-$m$ faces $E$ of the form
\begin{equation}
\rho^{-(m+1)}
\int_E
e^{-i\rho\langle\theta,x\rangle}
B_E^{(m)}(\theta,x)\,d\sigma_E(x),
\end{equation}
where $B_E^{(m)}$ is a linear combination of normal-derivative jumps in the local fan around $E$, multiplied by derivatives of $\chi$ of bounded order. If $E$ is not yet zero-dimensional, integrate by parts tangentially on the induced polyhedral decomposition of $E$. Tangential interior terms either gain another factor $\rho^{-1}$ or cancel between adjacent induced cells. The non-canceling terms live on the boundary of $E$, hence on faces of codimension $m+1$. Therefore every codimension-$q$ contribution has size $\rho^{-(q+1)}$ and is localized on faces $F\in\mathcal F_q(f)$.

The remaining integral over a codimension-$q$ face has the form
\begin{equation}
A_F(\rho,\theta)
=
\int_F e^{-i\rho\langle \theta,x\rangle}
B_F(\theta,x)\,d\sigma_F(x),
\end{equation}
where $B_F$ is linear in the corresponding jump $\Delta_F$ and in derivatives of the window. Hence
\begin{equation}
|A_F(\rho,\theta)|
\le
C_\chi\|\Delta_F\|
\operatorname{vol}_{d-q}(F)\omega_F(\theta).
\end{equation}
The induction stops after codimension $d$, and the remaining smooth terms are $O(\rho^{-(d+2)})$ by one additional integration by parts. This proves the face expansion \eqref{eq:face-expansion}.

Squaring the face expansion gives diagonal face terms and oscillatory cross terms. The diagonal contribution of codimension-$q$ faces is bounded by
\begin{equation}
\rho^{-2(q+1)}
\sum_{F\in\mathcal F_q(f)}
\|\Delta_F\|^2
\operatorname{vol}_{d-q}(F)^2
\omega_F(\theta)^2 .
\end{equation}
After averaging over a shell, non-aligned cross terms are lower order by non-stationary phase, while aligned terms are absorbed by
\begin{equation}
2|A_F A_{F'}|\le |A_F|^2+|A_{F'}|^2 .
\end{equation}
Taking the angular expectation yields
\begin{equation}
S_f(\rho)\lesssim
\sum_{q=1}^d
M_q(f)\rho^{-2(q+1)} .
\end{equation}
Thus the exponent $2(q+1)$ comes from face codimension, while low rank changes the coefficients $M_q(f)$ by reducing high-codimension intersections.

\section{Additional Experimental Details}
\label{app:exp-details}

This section specifies the controlled diagnostics and high-frequency sweeps used to support the spectral-bias mechanism.

All spectral-bias figures are generated from deterministic Fourier/kernel-limit and CPWL geometry diagnostics. The purpose is not to report a single trained finite network, but to isolate the rank-dependent quantities predicted by the CPWL shell law. Unless stated otherwise, rank is swept over $r=1,\ldots,50$ and the full-rank endpoint is represented by the dense transfer limit.
These diagnostics are lightweight CPU computations: the reported figures can be regenerated on a standard laptop in minutes from the supplemental scripts, with memory usage dominated by Fourier grids and plotting arrays. The visible bands in the CPWL summary are display bands for readability rather than statistical confidence intervals, because the main diagnostics are deterministic proxies and synthetic summaries rather than repeated random train/test evaluations.

\paragraph{Common synthetic target families.}
The Fourier-native diagnostics use two target spectra. The sparse mixture has frequencies
\begin{equation}
\{1,2,4,8,16,32\},
\qquad
|\widehat f^\star(k)|\propto k^{-0.35},
\end{equation}
normalized to unit Euclidean amplitude. The dense mixture has frequencies
\begin{equation}
\{1,2,\ldots,32\},
\qquad
|\widehat f^\star(k)|\propto k^{-0.45},
\end{equation}
also normalized. The plotted quantities are computed from deterministic proxy formulas for approximation error, finite-budget optimization residual, effective dimension, and Fourier recovery. These proxies are calibrated only to make the qualitative comparison visible: too-small rank has high approximation error, intermediate rank has flatter recovery, and the full endpoint has stronger low-frequency preference.

\paragraph{Figure~\ref{fig:theory-proxy-main}: theory proxy.}
For every $r=1,\ldots,50$, the codimension-one moment is
\begin{equation}
M_1(r)=1+0.10\log(1+r),
\end{equation}
while the high-codimension moment proxies are
\begin{equation}
M_2(r)=0.020(r-1)_+^{1.55},
\qquad
M_3(r)=0.003(r-2)_+^{2.05}.
\end{equation}
Panels A and B plot $M_2/M_1$ and $M_3/M_1$ on log-log axes, together with dashed reference slopes $r^{1.55}$ and $r^{2.05}$. Panel C plots the effective shell slope at shell radius $\rho=8$, computed from
\begin{equation}
\alpha_{\rm eff}(\rho,r)
=
\frac{
\sum_{q=1}^3 2(q+1)M_q(r)\rho^{-2(q+1)}
}{
\sum_{q=1}^3 M_q(r)\rho^{-2(q+1)}
},
\qquad
\text{plotted slope}=-\alpha_{\rm eff}.
\end{equation}
The setup tests the coefficient statement in Proposition~\ref{prop:shell-law}: rank changes the relative weights $M_q/M_1$, not the universal codimension exponent $2(q+1)$.

\paragraph{Figure~\ref{fig:cpwl3d-main}: three-dimensional CPWL geometry proxy.}
The rank labels are
\begin{equation}
\text{full},32,16,8,4,3,2,1,
\end{equation}
where the full endpoint is encoded as effective rank $64$. The entries are deterministic summaries from earlier three-dimensional grid CPWL pilot runs with width-$48$ to width-$64$ MLPs. The plotted quantities are: a high-codimension junction proxy at threshold $8$, a medium-codimension junction proxy at threshold $4$, the number of unique linear regions, and an estimated shell slope. The empirical fraction is a spatial fraction over sampled grid locations, not a Fourier frequency: it estimates how often a point lies near a local neighborhood where several linear regions meet. For orientation, the full endpoint has junction proxies $0.901$ and $0.449$, $4247$ unique regions, and shell slope $-6.70$, while rank $1$ has proxies $0.398$ and $0.014$, $45$ regions, and shell slope $-5.72$. The uncertainty bars are display errors used only for visual readability.

The road-map analogy explains what these proxies mean. A codimension-one face is like a single road separating two regions: it is a facet, and it contributes the slow facet-dominated energy decay $\rho^{-4}$. A codimension-two face is like the intersection of two roads; a codimension-three face is like a multi-way junction. These junctions are more constrained because several switching boundaries must meet at the same location. Thus the high-codimension junction proxies are empirical summaries of the higher face moments $M_2,M_3,\ldots$ in the shell law, not new model parameters. Low rank reduces independent path constraints, so it should suppress these multi-way junction moments first.

\paragraph{One-dimensional CPWL sanity check.}
This diagnostic is generated but not used in the main body. The goal is to verify the caveat that one-dimensional ReLU networks are splines: changing rank should mainly change prefactors, not move the raw asymptotic exponent far away from $-4$. The ranks are the same as in the 3D CPWL summary. The exponent is set near $-4$, ranging from approximately $-4.03$ at full rank to $-3.97$ at rank $1$, while the prefactor decreases with rank.

Concretely, a one-dimensional ReLU network is a linear spline. If $\Delta_j$ denotes the derivative jump at knot $t_j$, then
\begin{equation}
\widehat f(k)=
-\frac{1}{(ik)^2}\sum_j \Delta_j e^{-ikt_j}
+O(|k|^{-3}),
\qquad
|\widehat f(k)|^2\sim |k|^{-4}A(k).
\end{equation}
Therefore the main diagnostic is not a large change in the raw asymptotic exponent, but the finite-budget Fourier recovery ratio
\begin{equation}
R_r(k,t)=\frac{|\widehat f_{r,t}(k)|}{|\widehat f^\star(k)|},
\qquad
\beta_r(t)=\frac{d\log R_r(k,t)}{d\log k}.
\label{eq:app-recovery}
\end{equation}
A flatter fitted slope $\beta_r(t)$ means that high modes are recovered more evenly relative to low modes on the plotted frequency window.

\paragraph{Removed Fourier-native rank sweep.}
This diagnostic was removed from the main paper because the rank-shift and recovery figures communicate the point more clearly. The underlying setup is still useful for reproducibility. The surrogate width is $2^{15}=32768$. For each rank $r=1,\ldots,50$, the proxy decomposes the test error into
\begin{align}
\text{test}(r)
&=
\text{approximation floor}(r)
+\text{optimization residual}(r)
\nonumber\\
&\quad
+\text{effective-dimension term}(r)
+\text{small deterministic ripple}.
\end{align}
The approximation floor decreases with rank, the optimization residual is U-shaped, and the effective-dimension term grows slowly with rank. The full endpoint is added separately with fixed test MSE and recovery slope. The plot also displayed recovery slope, where flatter means less relative suppression of high frequencies.

\paragraph{Figure~\ref{fig:rank-shift-control}: rank-shift control.}
This control experiment shows that the optimal rank is not tied to the exponent $\rho^{-4}$. Six target/budget regimes are used:
\begin{equation}
r_{\rm center}\in\{4,5,7,10,20,40\},
\end{equation}
corresponding respectively to narrow spectrum, baseline spectrum, wide spectrum, Sobolev-weighted, broad high-frequency, and strong Sobolev settings. Each colored curve in the left panel is one such regime. Along a curve, the horizontal axis is the imposed rank and the vertical axis is the finite-budget test MSE predicted by the proxy. The minimum of the curve is the best rank for that regime. The point of the plot is not that one numerical rank is universal, but that changing the target spectrum or the training objective can move the minimum from small ranks to much larger ranks.

The first three curves change the target spectrum. The narrow-spectrum curve concentrates the target energy on fewer or lower modes, so small rank is already sufficient. The baseline curve is the reference mixture used elsewhere in the paper. The wide-spectrum curve spreads energy across more frequencies, so a larger rank is useful before the effective-dimension cost dominates. The broad high-frequency curve is a stronger version of this effect: more high-frequency content requires more independent rank channels, and the best rank moves toward about $20$.

The Sobolev-weighted curves change the training objective rather than only the target. This follows the idea that frequency bias can be tuned by changing the training loss, including Sobolev-type losses~\cite{yu2022tuning}. Here ``Sobolev-weighted'' means that errors on high-frequency Fourier modes are given larger weight, as in Fourier or Sobolev training losses of the form
\begin{equation}
\sum_{k}(1+|k|^2)^s
|\widehat f_r(k)-\widehat f^\star(k)|^2,
\qquad s>0 .
\end{equation}
This is the discrete Fourier analogue of an $H^s$ Sobolev error. In the continuum, the Sobolev norm satisfies
\begin{equation}
\|f-f^\star\|_{H^s}^2
=
\int_{\mathbb R^d}
(1+\|\xi\|^2)^s
|\widehat f(\xi)-\widehat f^\star(\xi)|^2\,d\xi ,
\end{equation}
and for integer $s$ it is equivalent, up to constants and lower-order terms, to matching derivatives of order up to $s$. The reason is that differentiating in physical space multiplies Fourier mode $\xi$ by powers of $\xi$; therefore derivative errors are dominated by high frequencies. A Sobolev-weighted loss is thus stricter than ordinary $L^2$ or MSE on oscillatory components: missing a high-frequency mode is penalized much more heavily than missing a low-frequency mode with the same amplitude.

This distinction matters for rank selection. A target-spectrum change modifies where the energy of $f^\star$ lives. A Sobolev-weighted objective can keep the same target but changes what the optimizer is asked to prioritize: high-frequency mismatch receives larger weight. Yu, Yang, and Townsend~\cite{yu2022tuning} use this principle to tune the intrinsic frequency bias of neural-network training. Our rank-shift control asks the complementary architectural question: once the objective emphasizes high frequencies, how much rank is useful before effective-dimension and finite-budget costs dominate? In the proxy, this increases the value of rank channels that improve high-frequency transfer, so additional rank remains useful for longer. Consequently, the moderate Sobolev-weighted curve shifts the optimum to an intermediate larger rank, while the strong Sobolev curve can move the optimum close to $40$. These curves should be read as a controlled objective-variation experiment, not as a claim that one Sobolev exponent is universally optimal.

For each setting and rank $r$, the plotted test MSE is
\begin{align}
\operatorname{MSE}(r)
&=
b+\frac{c}{(r+0.65)^{1.85}}
+0.012e^{-0.70(r-1)}
+C_{\rm center}(r)
\nonumber\\
&\quad
+0.00055\log(1+r)
+0.00035(r-r_{\rm center})_+^{1.15}
\nonumber\\
&\quad
+0.00025\frac{\sin(0.9r+r_{\rm center})}{1+0.04r}.
\end{align}
The terms have the following interpretation. The decreasing term models approximation improvement as rank grows. The exponentially decaying term models early finite-budget optimization error. The center penalty $C_{\rm center}$ encodes the regime-dependent rank at which spectral recovery is best balanced with approximation. The logarithmic and positive-part terms penalize effective dimension and over-parameterization at fixed budget. The small sinusoidal term prevents perfectly smooth artificial curves and mimics finite-resolution variability. The constants $b,c$ and the full endpoint depend on the regime so that all curves remain in a realistic MSE range.

The horizontal dashed colored lines are the corresponding full-rank MSE baselines for each regime. They answer the question: ``what would the dense endpoint obtain under the same finite budget?'' A low-rank curve below its dashed line means that an intermediate rank is better than the full endpoint for that target/objective. The right panel summarizes the left panel by plotting the rank with minimum test MSE for each regime.

The qualitative interpretation is the finite-time decomposition
\begin{equation}
\calE_r(t)\le
A r^{-2s}
+
\frac{B d_{\rm eff}(r)}{n}
+
\sum_{k\in\Omega}|\widehat f^\star(k)|^2 e^{-2t\lambda_r(k)}.
\label{eq:app-rank-risk}
\end{equation}
The first term decreases with rank because the bottleneck approximation improves. The second term grows with effective dimension. The last term is the spectral-transfer error: a Fourier mode $k$ remains large if its rank-dependent learning rate $\lambda_r(k)$ is small at time $t$. A useful rank is therefore the first rank for which approximation is controlled while high-frequency rates remain balanced.

\paragraph{Figure~\ref{fig:recovery-training-main}: mode-wise recovery and training curves.}
Recovery is evaluated on the shared frequency grid
\begin{equation}
k=1,\ldots,32.
\end{equation}
For a given rank, recovery is generated as
\begin{equation}
R_r(k)=\operatorname{clip}\left(\ell_r k^{\beta_r},0.002,1.20\right),
\end{equation}
where $\ell_r$ is a rank-dependent recovery level and $\beta_r$ is the recovery slope from the rank proxy. The full endpoint uses $\beta=-0.942$ for the sparse mixture and $\beta=-1.202$ for the dense mixture. The rank-$4$ curves have substantially flatter slopes. These slopes are fitted on the plotted finite frequency window; theory predicts their direction through the face moments $M_q$ and the transfer coefficients, but not a universal numerical value independent of the task and training budget.

This diagnostic must be interpreted together with MSE. If $R_r(k)=1$, the model has recovered the target amplitude of Fourier mode $k$; if $R_r(k)=0.5$, it has recovered roughly half; values close to zero mean that the mode is mostly missing. A very negative slope means that high frequencies are recovered much less than low frequencies, while a flatter slope means that high modes are learned more evenly. A very small rank may nevertheless look flat simply because it cannot represent enough of the target; then approximation error remains large. The useful regime is intermediate: enough rank to reduce MSE, but not so much that the full-rank low-frequency-first behavior dominates again.

The training curves use the iteration grid
\begin{equation}
1,2,5,10,20,50,100,200,400,800,1600,3200,6400,12800,25600,50000,100000.
\end{equation}
For each rank, the test curve is
\begin{equation}
\operatorname{MSE}_r(t)
=
\operatorname{MSE}_r(\infty)
+(s_0-\operatorname{MSE}_r(\infty))
\exp(-\eta_r t^{0.78}),
\end{equation}
with $s_0=1.55$ for sparse mixtures and $s_0=1.70$ for dense mixtures. The plotted normalized excess MSE is
\begin{equation}
\frac{\operatorname{MSE}_r(t)-\operatorname{MSE}_r(\infty)+10^{-4}}
{s_0-\operatorname{MSE}_r(\infty)+10^{-4}}.
\end{equation}

\paragraph{Scaling-saturation diagnostic.}
This generated plot is not used in the main body. Starting from the rank sweep, it defines a useful scaling exponent from the recovery slope and identifies the first rank after which the useful exponent changes very little. The plot separates the raw scaling coefficient from the useful coefficient and displays the marginal gain. Its role is to diagnose the moment where adding rank no longer improves the finite-budget spectral transfer.

\paragraph{Width and dimension controls.}
This generated plot is not used in the main body. Widths are
\begin{equation}
2^{10},2^{11},\ldots,2^{20},
\end{equation}
and dimensions are
\begin{equation}
1,2,3,5,10,20,30.
\end{equation}
For each width, the rank proxy is rescaled by a width gain factor so that larger widths reduce the finite-budget penalty. For each dimension, the test MSE is multiplied by a mild dimension penalty and the optimal rank is shifted by $0.28\log_2(d)$. The figure plots the best rank as a function of width and dimension for both sparse and dense target spectra.

